% Замените поясняющий текст собственными результатами.
% В шаблоне нет готового решения и обязательных внешних изображений.
\reportelement{Введение}

Цель работы: собираем свой образ приложения, уменьшаем его размер и добавляем проверку здоровья.~\ref{tab:parameters}

\begin{table}[H]
  \caption{Персональные параметры}\label{tab:parameters}
  \centering
  \begin{tabularx}{\textwidth}{|p{0.46\textwidth}|X|}
    \hline
    Параметр & Значение\\ \hline
    Номер ИСУ & \StudentISU\\ \hline
    Персональный порт & 5462\\ \hline
    Персональный префикс & karpov\\ \hline
    Окружение & linux/x86\_64\\ \hline
  \end{tabularx}
\end{table}

Результат хранится в ветке \texttt{lab-\LabNumber} форка:
\url{\RepositoryURL}. 

\chapter{Первый рабочий образ}

Собираем первый докер образ для приложения, запускаем контейнер и проверяем через команду curl (рис. 1):
\reportscreenshot{step1-4.png}{Результат проверки работы curl}{fig:proof1}

На рисунке~\ref{fig:proof1} видно что в ответе /me возвращается фамилия и в /secret виден флаг.


\chapter{Оптимизация образа}
Создадим новый файл Dockerfile.optimized, в котором настроим оптимизированную multistage сборку (рис. ~\ref{fig:proof2}): 

\reportscreenshot{step2multistage.png}{Multistage сборка}{fig:proof2}

Далее соберем образ и сравним размер со старым (рис. ~\ref{fig:proof3}):
\reportscreenshot{step2.png}{Сравнение образов}{fig:proof3}
Как видим, образ оказался меньше на 11 МБ. Связано это с тем, что промежуточные слои из предыдущего builder-этапа (временные данные pip) не попали в финальный образ, попали только нужные зависимости. Также отдельное копирование requirements позволяет ускорить последующие сборки, так как не нужно будет заново копировать и устанавливать зависимости из этого файла, если они не изменялись, возьмуться просто старые зависимости которые мы устанавливали при предыдущей сборке (докер не пересобирает слой если слой не был изменен).


\chapter{.dockerignore}

Соберем образ с .dockerignore (рис. ~\ref{fig:proof4}):
\reportscreenshot{step3.png}{образ с .dockerignore}{fig:proof4}
Без .dockerignore в образ мог попасть файл README.md и также могли бы попасть файлы окружения и файлы с секретами (паролями, токенами и т.д.), что могло бы привести к уязвимостям безопасности и утечке личной информации. Содержимое .dockerignore:
\reportscreenshot{dockerignore.png}{.dockerignore}{fig:proof10}


\chapter{Healthcheck - готовность к оркестрации}

Добавим проверку healthcheck и установим curl в финальный образ, а потом сравним размеры образов (рис. ~\ref{fig:proof5}):
\reportscreenshot{step4.png}{образ с curl и healthcheck}{fig:proof5}

Итоговый образ стал весить 178 МБ, больше на 14 МБ чем раньше. Можно было обойтись и без curl, использовав urllib из стандартной библиотеки Python. 
Далее запускаем контейнер и проверяем, что он healthy (рис. ~\ref{fig:proof6}):

\reportscreenshot{step4healthy.png}{образ с curl и healthcheck}{fig:proof6}


\chapter{Финальный артефакт и тегирование}

Финальный артефакт работы(рис. ~\ref{fig:proof7}):
\reportscreenshot{step5.png}{Артефакт}{fig:proof7}









\reportelement{Заключение}
В результате выполнения работы удалось собрать свой docker-образ, оптимизировать его через multi-stage сборку, настроить healthcheck. 

% Оставьте только использованные источники и цитируйте их в тексте.
\begin{thebibliography}{9}
  \bibitem{docker-docs}
  Docker Documentation: документация Docker.
  URL: \url{https://docs.docker.com/}
  (дата обращения: 09.11.2026).
\end{thebibliography}

